\section{Construcción de indicadores socioeconómicos}
\label{ch:indicadores}

\subsection{Principios de construcción}
\label{sec:principios_construccion}

Todo indicador se construye siguiendo el esquema:

\begin{equation}
I_j = \frac{N_j}{D_j}
\label{eq:indicador}
\end{equation}

donde $I_j$ es el indicador $j$, $N_j$ el numerador (conteo de la condición de vulnerabilidad) y $D_j$ el denominador (población de referencia). Cuando el indicador es una proporción, $I_j \in [0,1]$. Cuando es una tasa, se expresa por $10^3$ o $10^4$ habitantes según la magnitud.

Los principios rectores son:

\begin{enumerate}
    \item \textbf{Transparencia de fórmula:} Numerador, denominador y población de referencia explícitos en \texttt{especificacion\_indicadores.csv}.
    \item \textbf{Comparabilidad municipal:} Mismo denominador conceptual para todos los municipios.
    \item \textbf{Escala homogénea:} Proporciones $[0,1]$ para indicadores PCA; tasas estandarizadas para descripción.
    \item \textbf{Dirección de vulnerabilidad:} Valor alto $\Rightarrow$ mayor vulnerabilidad (Sección \ref{sec:direccion_vulnerabilidad}).
    \item \textbf{Tratamiento de NR:} Dos versiones (base total / base válida) para indicadores con categoría NR (Sección \ref{ch:tratamiento_faltantes}).
    \item \textbf{Denominador cero $\rightarrow$ NA + flag:} No se imputa, no se elimina el municipio.
\end{enumerate}

\subsection{Numeradores}
\label{sec:numeradores}

Los numeradores son conteos directos de las condiciones de vulnerabilidad extraídos de los cuadros censales seleccionados:

\begin{itemize}
    \item \texttt{1HM}: Hogares con personas menores de 15 años (Total, Cabecera, Centro poblado, Rural disperso).
    \item \texttt{2HM}: Hogares con personas de 60 años y más (Total, Cabecera, Centro poblado, Rural disperso).
    \item \texttt{5PM}: Personas de 10 años y más por actividad realizada la semana anterior (Ocupado, Desocupado, Inactivo, por sexo).
    \item \texttt{6V}: Viviendas ocupadas con personas presentes por material predominante de paredes exteriores.
    \item \texttt{8V}: Viviendas ocupadas con personas presentes por disponibilidad de servicios públicos (Energía, Gas, Acueducto, Alcantarillado, Recolección, Internet).
    \item \texttt{10V}: Personas censadas en hogares particulares por estrato de energía (1, 2, 3, 4, 5, 6, Sin estrato).
    \item \texttt{17PM}: Población de 5 años y más por asistencia escolar y sexo (pendiente verificación: ¿nivel educativo o asistencia?).
\end{itemize}

\subsection{Denominadores y poblaciones de referencia}
\label{sec:denominadores}

\begin{table}[htbp]
\centering
\caption{Denominadores por indicador}
\label{tab:denominadores}
\begin{adjustbox}{max width=\textwidth}
\begin{tabular}{@{}lll@{}}
\toprule
\textbf{Indicador} & \textbf{Denominador} & \textbf{Fuente denominador} \\
\midrule
Prop. hogares con menores & Total hogares (1HM Total) & 1HM \\
Prop. hogares con adultos 60+ & Total hogares (1HM Total) & 1HM \\
Proporción de categorías seleccionadas en actividad (10+) & Población 10+ años (1PM) & 1PM (grupos 10+) \\
Prop. viviendas sin energía & Viviendas ocupadas presentes (1V\_MPIO) & 1V\_MPIO \\
Prop. viviendas sin gas & Viviendas ocupadas presentes (1V\_MPIO) & 1V\_MPIO \\
Prop. viviendas sin recolección & Viviendas ocupadas presentes (1V\_MPIO) & 1V\_MPIO \\
Prop. viviendas sin internet & Viviendas ocupadas presentes (1V\_MPIO) & 1V\_MPIO \\
Prop. hogares estrato 1-2 & Personas en hogares particulares (10V Total) & 10V \\
Prop. población 5+ nivel educativo & Población 5+ años (17PM Total = 1PM 5+) & 17PM / 1PM \\
Prop. viviendas material paredes inadecuado & Viviendas ocupadas presentes (1V\_MPIO) & 1V\_MPIO \\
Prop. viviendas sin acueducto & Viviendas ocupadas presentes (1V\_MPIO) & 1V\_MPIO \\
Prop. viviendas sin alcantarillado & Viviendas ocupadas presentes (1V\_MPIO) & 1V\_MPIO \\
\bottomrule
\end{tabular}
\end{adjustbox}
\end{table}

\subsection{Verificación de población de referencia}
\label{sec:verif_pob_ref}

\begin{itemize}
    \item El total de \texttt{17PM} coincide exactamente con la población de 5+ años de \texttt{1PM} en los 1.122 municipios (validación ``Diferencias total 17PM vs población 5+ 1PM'': 0, PASS). Ver \texttt{Final\_Base.md}.
    \item \texttt{17PM} \textbf{no permite} filtrar categorías de nivel educativo restringidas a población de 10+ años; se calculan usando el total de 17PM y se documenta la discrepancia.
    \item \texttt{pob\_10mas} se calculó desde los grupos de 1PM de 10 años o más; la cobertura de 5PM se coteja contra ese valor (28 municipios con cobertura $<0,9$).
\end{itemize}

\subsection{Escala}
\label{sec:escala}

\begin{itemize}
    \item \textbf{Indicadores para PCA (provisionales):} Proporciones $[0,1]$ (7 indicadores).
    \item \textbf{Indicadores de descripción/solapamiento IPM:} Proporciones $[0,1]$ (7 indicadores).
    \item \textbf{No se usan conteos absolutos} en la matriz PCA (no son comparables entre municipios de distinto tamaño).
\end{itemize}

\subsection{Dirección de vulnerabilidad}
\label{sec:direccion_vulnerabilidad}

La decisión metodológica es orientar los indicadores para que:

\begin{equation}
\text{valor alto} \Rightarrow \text{mayor vulnerabilidad}
\end{equation}

\begin{table}[htbp]
\centering
\caption{Dirección de vulnerabilidad por indicador}
\label{tab:direccion_vulnerabilidad}
\begin{adjustbox}{max width=\textwidth}
\begin{tabular}{@{}lll@{}}
\toprule
\textbf{Indicador} & \textbf{Dirección} & \textbf{Nota} \\
\midrule
Prop. hogares con menores 15 & Alto = mayor vuln. & Carga demográfica infantil \\
Prop. hogares con adultos 60+ & \textbf{Pendiente} & Correlación negativa con privaciones IPM ($-0,06$ a $-0,63$); puede ser transición demográfica \\
Tasa desocupación 10+ & Alto = mayor vuln. & Privación laboral \\
Tasa inactividad 10+ & Alto = mayor vuln. & Privación laboral \\
Prop. viviendas sin energía & Alto = mayor vuln. & Privación servicios \\
Prop. viviendas sin gas & Alto = mayor vuln. & Privación servicios \\
Prop. viviendas sin recolección & Alto = mayor vuln. & Privación servicios \\
Prop. viviendas sin internet & Alto = mayor vuln. & Privación conectividad \\
Prop. hogares estrato 1-2 & Alto = mayor vuln. & Proxy ingreso bajo \\
Prop. población 5+ bajo nivel educativo & Alto = mayor vuln. & Privación capital humano (si 17PM es nivel educativo) \\
Prop. viviendas paredes inadecuadas & Alto = mayor vuln. & Privación vivienda (duplica IPM) \\
Prop. viviendas sin acueducto & Alto = mayor vuln. & Privación servicios (duplica IPM) \\
Prop. viviendas sin alcantarillado & Alto = mayor vuln. & Privación servicios (duplica IPM) \\
\bottomrule
\end{tabular}
\end{adjustbox}
\end{table}

\textbf{Nota crítica:} El indicador ``Hogares con adultos 60+'' muestra correlación \textbf{negativa} con casi todas las privaciones del IPM (Spearman $-0,06$ a $-0,63$): los municipios más envejecidos tienden a ser \textit{menos} vulnerables. Su dirección debe definirse conceptualmente antes del PCA (puede comportarse como variable de transición demográfica, no de privación). Ver \texttt{decisiones\_variables\_IVS\_CNPV2018.csv}, fila 5.

\subsection{Indicadores seleccionados (matriz candidata provisional)}
\label{sec:indicadores_seleccionados}

Siete indicadores conforman la \textbf{matriz candidata para diagnóstico PCA} (base total y base válida):

\begin{enumerate}
    \item \texttt{actividad\_desocupacion\_incapacidad\_hogar} (derivado de 5PM): agrupa búsqueda de trabajo, oficios del hogar, incapacidad permanente y otra situación. No es una tasa de desempleo y su interpretación como privación laboral es parcial; requiere revisión conceptual antes del PCA.
    \item \texttt{sin\_energia\_electrica} (derivado de 8V) -- alto = mayor vuln.
    \item \texttt{sin\_gas\_red\_publica} (derivado de 8V) -- alto = mayor vuln.
    \item \texttt{sin\_recoleccion\_basuras} (derivado de 8V) -- alto = mayor vuln.
    \item \texttt{sin\_internet} (derivado de 8V) -- alto = mayor vuln.
    \item \texttt{hogares\_con\_menores15} (derivado de 1HM): describe composición demográfica y no es una privación por sí misma.
    \item \texttt{estrato\_1\_2} (derivado de 10V): proxy territorial sujeto a limitaciones de cobertura rural; MSA $=0,414$ y dirección como vulnerabilidad requiere revisión.
\end{enumerate}

Estos 7 indicadores son \textbf{candidatos para diagnóstico, no una selección aprobada}. Cuatro representan servicios/conectividad, uno combina categorías heterogéneas de actividad, uno describe composición del hogar y uno usa estratificación de cobertura territorial desigual. Antes del PCA se debe justificar la cobertura conceptual de cada dominio y el peso implícito que recibiría.

\subsection{Indicadores descartados del PCA (descripción/solapamiento IPM)}
\label{sec:indicadores_descartados}

Siete indicadores se calculan para descripción y validación cruzada, pero \textbf{quedan fuera de la matriz PCA propuesta} por solapamiento directo con privaciones del IPM:

\begin{enumerate}
    \item Nivel educativo (17PM) -- duplica ``Bajo logro educativo'' / ``Analfabetismo'' IPM.
    \item Material inadecuado de paredes (6V) -- duplica ``Material inadecuado de paredes exteriores'' IPM.
    \item Sin acueducto (8V) -- duplica ``Sin acceso a fuente de agua mejorada'' IPM.
    \item Sin alcantarillado (8V) -- duplica ``Inadecuada eliminación de excretas'' IPM.
    \item Hacinamiento crítico (5V) -- duplica ``Hacinamiento crítico'' IPM.
    \item Trabajo infantil / informalidad / dependencia -- no construidos desde CNPV; ya en IPM.
    \item Hogares con adultos 60+ -- dirección pendiente; correlación negativa con IPM.
\end{enumerate}

\subsection{Indicadores reservados para validación}
\label{sec:indicadores_validacion}

\begin{itemize}
    \item \textbf{IPM Municipal (agregado):} 15 privaciones ponderadas. No entra en la matriz candidata. Se reserva para un contraste externo posterior, interpretado con cautela por el solapamiento conceptual entre algunas privaciones y variables censales. No se presenta aquí una validación empírica del índice.
    \item \textbf{IPM por dominio (Cabeceras, Centros poblados y Rural):} Solo para análisis de brechas urbano-rural futuro.
\end{itemize}

\subsection{Especificación técnica}
\label{sec:especificacion_tecnica}

La especificación vigente de las 17 definiciones operativas (fuente, cuadro, numerador, denominador, población de referencia, dirección esperada, propuesta de inclusión, notas) está en \texttt{resultados/base\_trabajo/especificacion\_indicadores.csv}. Este archivo es la \textbf{fuente única de verdad} para la construcción de indicadores y debe ser aprobado/modificado por el autor antes de ejecutar PCA definitivo.


